% Apresentacao enxuta do registro de um projeto; compartilhada com o exemplo.
\makeatletter
\renewcommand\section{%
  \@startsection{section}{1}{\z@}{0.8\baselineskip}{0.25\baselineskip}{\normalfont\bfseries}%
}
\makeatother
\begin{center}
  \includegraphics[width=4cm]{identidade/logo_horizontal.png}\par
  \medskip
  {\large\bfseries RELATÓRIO PARCIAL DO PROJETO}\par
  \small Registro de memória, conteúdo e evidências para o portfólio
\end{center}
\smallskip
\begingroup
\small
\renewcommand{\arraystretch}{1.12}
\begin{tabularx}{\textwidth}{@{}>{\bfseries}p{0.22\textwidth}X@{}}
Identificação & \registroid{} --- \registronome \\
Estudante & \registroautor \\
Curso e unidade & \registrocurso{} / \registrounidade \\
Período & \registroperiodo \\
Equipe & \registroequipe \\
Orientação & \registroorientacao \\
Anotação & \registrodata \\
GitHub & \registrorepositorio \\
Versão & \registroversao \\
\end{tabularx}
\endgroup
\par\medskip
{\small\noindent\textbf{Piloto para revisão:} síntese de até cinco páginas
por projeto, incluindo identificação, tabelas, figuras e fontes consultadas.
Materiais extensos ficam no acervo, ligados por identificador e versão.\par}
\medskip
